% Part: normal-modal-logic
% Chapter: sequent-calculus
% Section: introduction

\documentclass[../../../include/open-logic-section]{subfiles}

\begin{document}

\olfileid[id]{nml}{seq}{int}

\olsection{Pengantar}

Kalkulus sekuen bagi logika proposisional dapat diperluas dengan aturan
tambahan yang menangani~\iftag{prvBox}{$\Box$\iftag{prvDiamond}{ dan~}{}}{}%
\iftag{prvDiamond}{$\Diamond$}{}. Sebagai contoh, bagi~$\Log{K}$, kita
mempunyai \Log{LK} ditambah:
\[
  \iftag{prvBox}{\iftag{prvDiamond}{% <> and [] primitive
    \Axiom$\Gamma \fCenter \Delta, !A$
    \RightLabel{$\Box$}
    \UnaryInf$\Box\Gamma \fCenter \Diamond\Delta, \Box!A$
    \DisplayProof
    \qquad
    \Axiom$!A, \Gamma \fCenter \Delta$
    \RightLabel{$\Diamond$}
    \UnaryInf$\Diamond!A, \Box\Gamma \fCenter \Diamond\Delta$
    \DisplayProof
}{% only Box primitive
\Axiom$\Gamma \fCenter !A$
\RightLabel{$\Box$}
\UnaryInf$\Box\Gamma \fCenter \Box!A$
\DisplayProof
}}{% only <> primitive
  \Axiom$!A \fCenter \Delta$
  \RightLabel{$\Diamond$}
  \UnaryInf$\Diamond!A \fCenter \Diamond\Delta$
  \DisplayProof
}
\]
Bagi perluasan~$\Log{K}$, aturan tambahan lain juga harus ditambahkan.

Tidak setiap logika modal mempunyai kalkulus sekuen semacam itu. Bahkan
$\Log{S5}$, yang secara semantik sederhana (logika itu dapat didefinisikan
tanpa menggunakan relasi aksesibilitas sama sekali), belum diketahui
mempunyai kalkulus sekuen yang diperoleh dari~$\Log{LK}$ dan lengkap tanpa
aturan~\Cut. Akan tetapi, logika itu mempunyai kalkulus \emph{hipersekuen}
lengkap yang bebas potongan.

\end{document}

